\documentclass[11pt,a4paper]{article}
\usepackage[T1]{fontenc}
\usepackage[utf8]{inputenc}
\usepackage{lmodern}
\usepackage{amsmath,amssymb}
\usepackage[margin=25mm]{geometry}
\usepackage[hidelinks]{hyperref}
\hypersetup{pdftitle={Halving goes to the ultraviolet; confinement lives in the infrared},pdfauthor={navstokgap research working notes}}
\newcommand{\tightlist}{\setlength{\itemsep}{0pt}\setlength{\parskip}{0pt}}
\setlength{\emergencystretch}{3em}
\setcounter{secnumdepth}{0}
\begin{document}
\hypertarget{halving-goes-to-the-ultraviolet-confinement-lives-in-the-infrared}{%
\section{Halving goes to the ultraviolet; confinement lives in the
infrared}\label{halving-goes-to-the-ultraviolet-confinement-lives-in-the-infrared}}

\textbf{Didactic note, 2026-09-28 (Claude, prompted by the user's
question).} Confinement and the mass gap are infrared properties, while
halving the lattice moves toward the ultraviolet. The refinement
programme of the \href{halving-atlas.md}{atlas} nevertheless bears on
the gap, for five reasons, each stated with its source:

\begin{enumerate}
\def\labelenumi{\arabic{enumi}.}
\tightlist
\item
  \textbf{One map, two directions.} A halving adds modes at the
  ultraviolet end; integrating them out again is one Wilson
  renormalization step toward the infrared. Refinement and
  coarse-graining are the same pushforward, Proposition 1 of the
  \href{series-parallel-gauge-refinement.md}{series/parallel note}, read
  in opposite directions.
\item
  \textbf{The infrared is the invariant of a refinement.} A continuum
  limit holds the physical scales fixed. In lattice units every infrared
  quantity shrinks at each halving: \(am\) and \(a\Lambda\) halve,
  \(a^2\sigma\) quarters. The question a halving asks about the infrared
  is whether these stay constant in physical units.
\item
  \textbf{The exponential ladder selects what survives.} Along the
  one-loop trajectory a non-perturbative term \(e^{-c/t}\), with
  \(t=\hbar g^2(a)\), equals \((a\Lambda)^{2b_0c}\), so a halving
  multiplies it by \(2^{-2b_0c}\). The generated gap,
  \(am\propto e^{-1/(2b_0\hbar g^2)}\), is the term with \(2b_0c=1\): it
  halves exactly when the lattice halves, and so survives in physical
  units. Every term with \(2b_0c>1\) dies.
\item
  \textbf{A growing ultraviolet end and a fixed infrared end.} From
  spacing \(a\) to the correlation length \(\xi\) there are
  \(\log_2(\xi/a)\) halvings. Refining adds steps only at the
  ultraviolet end, where \(t\) is small; the steps where \(t\) is of
  order one, around \(\beta_W\simeq6\) for \(SU(3)\), are the same
  finite set for every \(a\). Confinement is decided there once;
  refinement requires the growing number of ultraviolet steps to be
  controlled uniformly.
\item
  \textbf{Refinement exposes lattice artifacts.} Compact \(U(1)\) has a
  lattice gap that refinement at fixed coupling removes. A property that
  survives every halving is a property of the continuum theory.
\end{enumerate}

\hypertarget{one-map-two-directions}{%
\subsection{1. One map, two directions}\label{one-map-two-directions}}

Proposition 1 of the series/parallel note writes the law of the coarse
links as the pushforward of the fine heat-kernel measure under the
product of half-links. Read from coarse to fine it is a refinement: the
fine measure is a consistent extension of the coarse one, with new modes
at the shortest scale. Read from fine to coarse it is a blocking step:
the mid-plane integral \(\Psi\) is the effective action generated by the
modes of wavelength between \(a\) and \(2a\). The same factor \(\Psi\),
and the same parallel defect, carry the renormalization in both
readings. What the atlas calls ``what survives refinement'' is, in the
other reading, ``what is left of the ultraviolet after coarse-graining
down to the scale of interest''.

\hypertarget{the-infrared-is-what-a-refinement-holds-fixed}{%
\subsection{2. The infrared is what a refinement holds
fixed}\label{the-infrared-is-what-a-refinement-holds-fixed}}

A lattice at spacing \(a\) has three scales: \(a\), the correlation
length \(\xi=\hbar/(mc)\) and the box \(L\). A continuum limit sends
\(a/\xi\to0\) at fixed \(\xi/L\) (or with \(\xi/L\to0\) afterwards). In
\(1+3\), one-loop running ties the bare coupling to the spacing through
\(a\Lambda=e^{-1/(2b_0t)}\), \(t=\hbar g^2(a)\), \(b_0=11N/(48\pi^2)\),
so that \(\Lambda\) does not move while \(a\) does. The infrared
quantities are ratios to \(\Lambda\): \(m/\Lambda\),
\(\sigma/\Lambda^2\). A halving never computes them; it tests whether
they remain what they were. This is the precise sense in which the
programme's gap statement (AGENTS.md: which operator, which limits in
which order, what the gap is measured in) is a statement about all
halvings at once.

In \(1+2\) the relation is simpler: the heat time \(t=\lambda_3a\)
itself halves at every step, the ultraviolet steps are close to
Gaussian, and the infrared scale is \(\hbar c\lambda_3\) (the cell-7
target \(E=C_3\hbar c\lambda_3\) of the atlas), fixed by the
dimensionful coupling.

\hypertarget{the-exponential-ladder}{%
\subsection{3. The exponential ladder}\label{the-exponential-ladder}}

At one loop, \(e^{-c/t}=(a\Lambda)^{2b_0c}\), and \(2b_0c\) is the
scaling dimension of the term (atlas §4, ``The exponential ladder''). A
mass has dimension one, so a physical gap is the term that halves with
the lattice: \(am\propto(a\Lambda)^1\), the member of the ladder with
\(2b_0c=1\). The corrections a single step leaves sit far higher: the
bridge image terms at \(2b_0c=11N/3\) (roots) and \(11N\) (the \(SU(3)\)
centre images of a trisection, \href{sun-midpoint-centre.md}{centre
note}), beyond the per-volume threshold \(2b_0c=4\) of the
\href{zero-spacing-any-action.md}{zero-spacing note}. So every
ultraviolet step contains the infrared scale as a tiny non-perturbative
term of exactly the size that survives, and nothing else of its kind
survives with it. This is dimensional transmutation read step by step.
The term is invisible to perturbation theory, since
\(e^{-1/(2b_0\hbar g^2)}\) has a vanishing Taylor series in \(g^2\); it
is visible to the exact step, which is why the atlas insists on exact
halvings.

\hypertarget{a-growing-ultraviolet-end-and-a-fixed-infrared-end}{%
\subsection{4. A growing ultraviolet end and a fixed infrared
end}\label{a-growing-ultraviolet-end-and-a-fixed-infrared-end}}

The trajectory from spacing \(a\) to the scale \(\xi\) has
\(\log_2(\xi/a)\) halvings. As \(a\to0\), every new halving is added at
the ultraviolet end, where \(t=\hbar g^2(a)\to0\) and the atlas cells
1--6 give one-step control. The infrared end is the handful of order-one
steps where \(t\) is of order one; for \(SU(3)\) with the Wilson action
these sit around \(\beta_W\simeq6\), a few steps above the correlation
length (the \href{confinement-scale-bands.md}{confinement-scale bands
note}, with published numbers). That end is the same finite set for
every \(a\). Its analysis is the
\href{intermediate-region-finite-verification.md}{finite verification}
of the intermediate region; the uniform control of an unbounded number
of ultraviolet steps feeding it is what Jaffe and Witten name as the
missing idea. The division of labour in the atlas follows: cells 1--6
are the ultraviolet steps, cell 7 is the infrared end.

\hypertarget{where-confinement-can-be-lost}{%
\subsection{5. Where confinement can be
lost}\label{where-confinement-can-be-lost}}

The compact \(U(1)\) theory in \(1+2\) is the proved example. Its
lattice gap comes from monopoles; at fixed coupling the refinement
removes it and the continuum limit is the free photon (Gross), while
along the Göpfert--Mack trajectory the gap survives only with a
diverging monopole density per physical volume (atlas §4; the
\href{villain-monopole-refinement.md}{monopole note}). The ultraviolet
step is where a lattice gap is exposed as a lattice artifact. In
\(SU(3)\) the coupling must run to zero logarithmically (cell 4), and
the gap is expected to survive because it is tied to \(\Lambda\), which
the running holds fixed.

\hypertarget{newtons-parallel}{%
\subsection{6. Newton's parallel}\label{newtons-parallel}}

In Newton's cell the sagitta and the Galileo action vanish under
refinement, \(K_\tau\propto\tau^3\)
(\href{cut-measure-newton.md}{cut-measure note}), while the force
survives as the ratio those vanishing quantities define. On the lattice
\(am\to0\) while \(m/\Lambda\) survives. In both cases refinement never
computes the surviving quantity directly: it certifies that a ratio
stays fixed while every per-cell quantity goes to zero. What refinement
cannot remove in either case is the supplied action floor \(\hbar\)
(atlas §4, ``The action floor survives what removes the mass gap'').

\hypertarget{consequence-for-state}{%
\subsection{7. Consequence for STATE}\label{consequence-for-state}}

None; this note explains the programme's logic for the joint paper and
for readers of the atlas.
\end{document}
